\makeatletter
\def\input@path{{./}{../}{../../}{../../../}{../../../../}{preamble/}{../preamble/}{../../preamble/}{../../../preamble/}{../../../../book2026-2027/preamble/}}
\makeatother
\input{preamble_exams.tex}
\input{graphs/highlights.tex}
% \GTenHighlightsfalse % Uncomment to disable instructional highlights.
% Population/model, sample evidence, design threshold, and result roles.
\newcommand{\Pop}[1]{\HighlightStroke{GTenHLParameter}{#1}}
\newcommand{\Obs}[1]{\HighlightStroke{GTenHLValue}{#1}}
\newcommand{\Design}[1]{\HighlightStroke{GTenHLThird}{#1}}
\newcommand{\Result}[1]{\HighlightStroke{GTenHLOutcome}{#1}}

\title{Grade 11 Global Economics\\Term 3 --- Lesson 1\\Home Practice C5: Tests for a population mean}
\author{}\date{}
\hypersetup{pdftitle={Grade 11 Global Economics - Term 3 - Lesson 1 - C5},pdfsubject={Tests for a population mean}}
\begin{document}
\maketitle
\setcounter{tocdepth}{1}
\setcounter{secnumdepth}{2}
\phantomsection\label{toc}
\tableofcontents
\section*{Evaluated criteria}
\begin{description}
\item[\textbf{C5}] Conducts hypothesis tests for a population mean using the Z-test when the variance is known and the t-test when the variance is unknown, selecting the appropriate statistic and critical value.
\end{description}
\textbf{How to practise.} Choose the model from what is known about population variance. Keep statistic, critical value, p-value, and level separate. Attempt each task before opening its solution link. Use the return arrows to revisit the problem, group, or contents. All datasets and scenarios in this activity are illustrative.
\clearpage
\section{Introduction}\label{sec:introduction}
\subsection*{Highlight legend}
\phantomsection\label{tab:highlight-legend}
\begin{tabularx}{\linewidth}{@{}>{\raggedright\arraybackslash}p{.22\linewidth}>{\raggedright\arraybackslash}p{.25\linewidth}>{\raggedright\arraybackslash}X@{}}
\toprule
\textbf{Symbol or notation}&\textbf{Concept}&\textbf{Brief meaning}\\\midrule
$\Pop{\mu},\Pop{\mu_0}$ & \Pop{Mean and benchmark} & Population value and reference. \\
$\Pop{\sigma},\Obs{s}$ & \Obs{Spread information} & Known versus estimated SD. \\
$\Obs{z},\Obs{t}$ & \Obs{Statistic} & Standardized evidence. \\
$\Design{\alpha},\Design{c}$ & \Design{Decision design} & Level and critical boundary. \\
\bottomrule
\end{tabularx}
\subsection{Introductory example 1: Known invoice variability}\label{ex:intro-one}
For $n=36$ invoices, compare $\bar x=79$ USD with 75 USD. Population variance is genuinely known as 144 USD$^2$, so $\Pop{\sigma}=\Pop{12}$ USD. An increase question at level 0.05 uses Z if independent sampling and a normal sampling distribution are justified.
\subsection{Introductory example 2: Estimated fee variability}\label{ex:intro-two}
For $n=25$ fees, compare $\bar x=18$ USD with 20 USD. Only sample $\Obs{s}=\Obs{5}$ USD is known. A decrease question uses t with $df=24$, subject to suitable sampling and shape conditions.
\section{Notation Reference and Practice Procedure}\label{sec:notation-and-procedure}
\subsection{Notation reference}\label{tab:notation-reference}
\begin{tabularx}{\linewidth}{@{}>{\raggedright\arraybackslash}p{.16\linewidth}>{\raggedright\arraybackslash}p{.18\linewidth}>{\raggedright\arraybackslash}p{.28\linewidth}>{\raggedright\arraybackslash}X@{}}
\toprule
\textbf{Symbol or notation}&\textbf{Concept}&\textbf{Definition}&\textbf{Meaning in the introductory examples}\\\midrule
$\Pop{\mu},\Pop{\mu_0}$ & \Pop{Mean and benchmark} & Both have the measurement units; hypotheses concern the population mean. & \textbullet\ \hyperref[ex:intro-one]{Example 1}: invoice mean and 75 USD.\par\textbullet\ \hyperref[ex:intro-two]{Example 2}: fee mean and 20 USD. \\
$\Pop{\sigma},\Obs{s}$ & \Obs{Spread information} & $\sigma=\sqrt{\sigma^2}$ is population SD; $s$ estimates it from the sample. & \textbullet\ \hyperref[ex:intro-one]{Example 1}: known $\sigma=12$.\par\textbullet\ \hyperref[ex:intro-two]{Example 2}: estimated $s=5$. \\
$\Obs{z},\Obs{t}$ & \Obs{Statistic} & $z=(\bar x-\mu_0)/(\sigma/\sqrt n)$; $t=(\bar x-\mu_0)/(s/\sqrt n)$, $df=n-1$. & \textbullet\ \hyperref[ex:intro-one]{Example 1}: Z model.\par\textbullet\ \hyperref[ex:intro-two]{Example 2}: t model, $df=24$. \\
$\Design{\alpha},\Design{c}$ & \Design{Decision design} & Critical values come from the null distribution and tail. A p-value is the probability of equally or more extreme evidence under that model. & \textbullet\ \hyperref[ex:intro-one]{Example 1}: right-tail level 0.05.\par\textbullet\ \hyperref[ex:intro-two]{Example 2}: left-tail level 0.05. \\
\bottomrule
\end{tabularx}
\subsection{Practice procedure}
\begin{enumerate}
\item Define the population mean, units, benchmark, claim direction, and planned level.
\item Check random independent sampling. Small samples require an approximately normal population without serious outliers; large samples still require care with heavy tails.
\item Use Z for known population variance; otherwise use estimated $s$ and t with $df=n-1$. Large $n$ alone does not make variance known.
\item Calculate standard error and statistic. Use a critical value and p-value from the same null distribution and tail.
\item Reject $H_0$ when p-value $\le\alpha$ or the statistic is in the rejection region; otherwise fail to reject. Interpret practical size separately.
\end{enumerate}
\section{General Worked Problems}\label{sec:general-worked}
\SectionProblemsTableOfContents{\SectionProblemLink{prob:known}{Known-variance invoice test} \\
\SectionProblemLink{prob:unknown}{Unknown-variance fee test} \\
\SectionProblemLink{prob:large}{A large sample with unknown variance}}
\subsection{Known-variance invoice test}\label{prob:known}
A random independent sample of 36 invoices from an approximately normal population has $\bar x=79$ USD. Known population variance is 144 USD$^2$. Test an increase above 75 USD at $\alpha=0.05$. Supplied: right critical 1.645; $P(Z\ge2)=0.0228$.
\QuestionsTasks[sec:general-worked]{{5}/{sol:known}/{Conduct the test and identify the four decision quantities.}}
\subsubsection*{Worked solution}\label{sol:known}
Population mean invoice cost is $\Pop{\mu}$; $H_0:\mu=75$, $H_a:\mu>75$. Independent random sampling and approximate normality are supplied. Known variance gives $\sigma=12$ USD and a Z model:
\[
SE=12/\sqrt{36}=2,\qquad z=\frac{\Obs{79}-\Pop{75}}{\Pop{12}/\sqrt{\Obs{36}}}=2.
\]
Statistic 2 exceeds critical $\Design{1.645}$; p-value 0.0228 is below level $\Design{0.05}$. Therefore \Result{reject $H_0$}. Evidence supports a population mean above 75 USD. The estimate is 4 USD per invoice; purchase volume and budget determine its economic importance.
\secInicio[prob:known][sec:general-worked]
\subsection{Unknown-variance fee test}\label{prob:unknown}
An independent random sample of 25 fees has $\bar x=18$ USD and $s=5$ USD. Population variance is unknown. The population is approximately normal without serious outliers. Test a decrease below 20 USD at $\alpha=0.05$. For $df=24$, supplied left critical is $-1.711$ and $P(T_{24}\le-2)\approx0.0285$.
\QuestionsTasks[sec:general-worked]{{5}/{sol:unknown}/{Conduct the test and explain why a t distribution is needed.}}
\subsubsection*{Worked solution}\label{sol:unknown}
Let $\Pop{\mu}$ be population mean fee: $H_0:\mu=20$, $H_a:\mu<20$. Conditions are supplied. Unknown population spread requires t with $df=25-1=24$:
\[
SE=5/\sqrt{25}=1,\qquad t=\frac{\Obs{18}-\Pop{20}}{\Obs{5}/\sqrt{\Obs{25}}}=-2.
\]
Since $-2<-1.711$ and p-value $0.0285<0.05$, \Result{reject $H_0$}. Evidence supports a mean below 20 USD; the sample estimates a 2 USD reduction. Estimating spread introduces uncertainty represented by the t distribution.
\secInicio[prob:unknown][sec:general-worked]
\subsection{A large sample with unknown variance}\label{prob:large}
An independent random sample of 100 export invoices has $\bar x=101$ USD and $s=10$ USD. Population variance is unknown; no heavy tails or influential outliers are present. Test any difference from 100 USD at $\alpha=0.05$. Supplied for $df=99$: two-tail critical values $\pm1.984$ and $P(|T_{99}|\ge1)\approx0.3197$.
\QuestionsTasks[sec:general-worked]{{5}/{sol:large}/{Select and conduct the test, explaining the role of sample size.}}
\subsubsection*{Worked solution}\label{sol:large}
For population mean $\mu$, use $H_0:\mu=100$, $H_a:\mu\ne100$. Independent random sampling and a large sample without severe shape problems support the approximation. Unknown $\sigma$ requires t, not Z solely because $n=100$. Here $df=99$, $SE=10/\sqrt{100}=1$, and $t=(101-100)/1=1$. Since $|1|<1.984$ and $0.3197>0.05$, \Result{fail to reject $H_0$}. Evidence is insufficient for a population mean difference. This does not prove equality; the estimated increase is 1 USD with uncertainty.
\secInicio[prob:large][sec:general-worked]
\section{Advanced Problems}\label{sec:advanced}
\SectionProblemsTableOfContents{\SectionProblemLink{prob:compare}{Same number, different information}}
\subsection{Same number, different information}\label{prob:compare}
For an independent random sample from a normal population, $n=16$, $\bar x=79$ USD, benchmark 75 USD, and the preplanned alternative is two-sided at $\alpha=0.05$. In A, population $\sigma=8$ USD is known. In B, only sample $s=8$ USD is known. Supplied: Z critical $\pm1.960$, two-sided p at 2 is 0.0455; $df=15$ t critical $\pm2.131$, two-sided p at 2 is 0.0639.
\QuestionsTasks[sec:advanced]{{5}/{sol:compare}/{Conduct both tests and explain different decisions and practical significance.}}
\subsubsection*{Worked solution}\label{sol:compare}
Both use $H_0:\mu=75$, $H_a:\mu\ne75$, with appropriate independent normal sampling. Both have numerical $SE=8/\sqrt{16}=2$ and statistic $(79-75)/2=2$. In A, known variance gives Z: $2>1.960$ and $0.0455<0.05$, so \Result{reject $H_0$}. In B, estimated variance gives t with $df=15$: $2<2.131$ and $0.0639>0.05$, so \Result{fail to reject $H_0$}. Heavier t tails reflect uncertainty in estimating spread. The 4 USD estimated difference is the same in both cases; differing decisions do not change economic size. Compare it with business costs before acting.
\secInicio[prob:compare][sec:advanced]
\end{document}

